% X^∞ -- Postmoral and Emotionless
% Mathematical Foundations of Ethical Governance as a Self-Reinforcing System
% Working Paper
% License: CC BY 4.0
% Author: The Auctor

\documentclass[12pt]{scrartcl}
\usepackage[utf8]{inputenc}
\usepackage[T1]{fontenc}
\usepackage[english,ngerman]{babel}
\usepackage{microtype}
\usepackage[a4paper,margin=2.5cm]{geometry}
\usepackage{graphicx}
\usepackage{xcolor}
\usepackage{amsmath,amssymb}
\usepackage{csquotes}
\usepackage{enumitem}
\usepackage{booktabs}
\usepackage{fancyhdr}
\usepackage{titlesec}
\usepackage{caption}
\usepackage{makecell}

% Font
\usepackage[sfdefault]{FiraSans}
\usepackage{inconsolata}

% Section formatting
\titleformat{\section}[block]{\Large\bfseries\sffamily}{\thesection}{1em}{}
\titleformat{\subsection}[block]{\large\bfseries\sffamily}{\thesubsection}{1em}{}
\titleformat{\subsubsection}[block]{\normalsize\bfseries\sffamily}{\thesubsubsection}{1em}{}

% Improve line breaking
\emergencystretch 3em

% Load hyperref last to avoid conflicts
\usepackage[unicode=true,pdfencoding=auto]{hyperref}

% Title and Author
\title{\texorpdfstring{Physics of Symbiotic Being X$^\infty$}{Physics of Symbiotic Being X-Infinity}\\Mathematical Foundations of Ethical Governance as a Self-Reinforcing System}
\author{The Auctor \\
\texttt{x\_to\_the\_power\_of\_infinity@protonmail.com} \\
\texttt{https://mastodon.social/@The\_Auctor} \\
\texttt{https://www.linkedin.com/in/der-auctor-b12375362/} \\
\href{https://zenodo.org/communities/xtothepowerofinfinity/}{zenodo.org} \\
\href{https://github.com/Xtothepowerofinfinity/Philosophie_der_Verantwortung}{GitHub: Xtothepowerofinfinity}
}
\date{June 1, 2025}

\begin{document}
\selectlanguage{english}
\maketitle
X$^\infty$ is neither a utopia nor a mere model – it is the Theory of Everything, unifying physical, social, and ethical systems under the universal principle \emph{actio $\Rightarrow$ reactio}. Just as gravitation describes the motion of planets, X$^\infty$ explains the dynamics of all interactions, free from the distortions of current power systems. Its mechanisms are already observable in reality.

\clearpage

\begin{abstract}
X$^\infty$ is a radical responsibility system based on the universal natural law \emph{actio $\Rightarrow$ reactio}. Emerging from the observation of everyday interactions, it formalizes the immediate feedback of actions and their consequences. Existing systems often suffer from power imbalances that deflect reactions and externalize responsibility. X$^\infty$ restores this natural law by ensuring the traceability of every effect to its source through a minimal set of fundamental mechanisms – in particular, the direct modification of action capacity ($\text{Cap}_{\text{Potential}}$) by effect ($\Delta$), a reciprocally weighted feedback ($w_{E'}$), a dynamic system efficiency quotient ($L$), and a decentralized petition system. This working paper presents the mathematical foundations of X$^\infty$ and serves for public discussion and validation.
\end{abstract}

\clearpage
\tableofcontents
\clearpage

\section{Fundamental Principle and Origin of \texorpdfstring{X$^\infty$}{X-Infinity}}
\label{sec:Grundprinzip}

X$^\infty$ is based on a universal natural law: Every action generates a reaction (\emph{actio $\Rightarrow$ reactio}). This principle, known from Newton’s third law of physics, applies equally to social systems. Every action of an entity produces effects that entail consequences. In traditional systems, however, these reactions are often deflected by power imbalances, so that the consequences frequently affect the weaker, the general public, or the environment rather than the perpetrators. X$^\infty$ restores the natural order by ensuring that every effect is traced back to its source. The origin of this model lies in the observation of a living community with children and pets, where actions and their consequences are immediately tangible. These dynamics revealed a fundamental principle: Actions generate effects that cannot be ignored but must be borne. Power imbalances distort this principle. X$^\infty$ formalizes this insight into a system that ensures the traceability of reactions and neutralizes power imbalances to create a natural, fair, and stable order. This natural law permeates all mechanisms of X$^\infty$:
\begin{itemize}
    \item \textbf{Feedback and $\text{Cap}_{\text{Potential}}$-Dynamics}: The action capacity ($\text{Cap}_{\text{Potential}}$) of an entity is directly modified by the effect ($\Delta$) of its actions. Positive effects can increase $\text{Cap}_{\text{Potential}}$, negative effects reduce it. This change is modulated by a dynamic system efficiency quotient ($L$), which accounts for systemic inefficiencies as “waste heat.” The feedback from other entities, particularly the relatively weaker ones (weighted by \( w_{E'} = \frac{1}{\max(1, \text{Cap}_{\text{Potential}}(E'))} \)), determines the registered effect on a subjective basis, as any evaluation of effect can only be subjective. A so-called “objective” evaluation leads to an objectifying bias that ignores the individual perspectives of those affected and can reinforce power imbalances.
    \item \textbf{Petitions}: Individual entities can articulate needs through petitions, whether systemic needs or the need to contribute their capabilities. Each petition originates from exactly one entity. Petitions can be directed to the broker to assign suitable tasks or directly articulate systemic needs. Their priority is derived from the feedback weight $w_{E'}$ of the petitioning entity. The adoption and execution of tasks occur exclusively through petitions from entities that voluntarily agree to perform them, thereby correcting power imbalances in task allocation.
    \item \textbf{Responsibility Conservation}: The totality of systemic necessities and the effects applied to them strives for a balance, analogous to a physical conservation law. This principle finds its ultimate expression in the function of the Lowest of the Low (UdU).
    \item \textbf{UdU (Lowest of the Low)}: Guarantees that no effect remains unborne and ensures fundamental systemic responsibility even in extreme situations, thus securing traceability and the principle of responsibility conservation.
\end{itemize}

The principle \emph{actio $\Rightarrow$ reactio} is postmoral: It does not evaluate intent but measures effect. X$^\infty$ is thus not a moral system in the traditional sense but a structure that mathematically describes and restores the natural law of social interactions.

\subsection{Objectives of the System}
The objectives remain fundamental but arise from the core structure:
\begin{itemize}
    \item \textbf{Fairness through Effect}: Responsibility and action opportunity are determined not by status or intent but by actual, subjectively fed-back consequences.
    \item \textbf{Protection of the Weakest (Emergent)}: Relatively weaker actors (with lower $\text{Cap}_{\text{Potential}}$ in the relevant context) receive a structurally stronger voice in the subjective evaluation of effects through reciprocal feedback weighting ($w_{E'}$).
    \item \textbf{Anti-Speciesism}: Humans, non-humans, and the environment are considered entities whose effects and affectedness are treated according to the same principles.
    \item \textbf{Self-Reinforcement and Self-Regulation}: The system learns and adapts through the direct effect of subjective feedback on $\text{Cap}_{\text{Potential}}$ and the dynamic system quotient $L$, without external normative interventions.
\end{itemize}

\subsection{\texorpdfstring{X$^\infty$}{X-Infinity} in Everyday Reality}
\label{sec:Alltagsrealitaet}

The mechanisms of X$^\infty$ are not abstract but observable in everyday interactions. The following examples illustrate how X$^\infty$ operates in real contexts:
\begin{itemize}
    \item \textbf{Family}: A child does not tidy up their toys, leading to chaos in the household. This affects all housemates who experience the disorder. The natural consequences – such as the need to tidy up themselves or the impairment of the shared living space – are the effect that directly traces back to the responsible entity (the child). X$^\infty$ formalizes this immediate feedback.
    \item \textbf{Company}: A poor decision by a decision-maker leads to a customer leaving the company. The direct feedback from the customer (e.g., through complaints or lost revenue) immediately affects the decision-maker, whose $\text{Cap}_{\text{Potential}}$ in the company context decreases. X$^\infty$ ensures that responsibility is not shifted to others.
    \item \textbf{Environment}: Pollution by an entity (e.g., a company) harms the environment and the entities affected by it (e.g., residents, animals). The evaluation by those directly experiencing the pollution matters most and determines the effect $\Delta$. X$^\infty$ ensures that this subjective experience shapes the consequences for the perpetrator.
    \item \textbf{Indigenous Wisdom}: Traditional cultures, such as indigenous communities that have existed stably for millennia, intuitively demonstrate the principles of X$^\infty$. Their harmonious coexistence with nature shows how actions and consequences are aligned with physical laws, promoting sustainable systems.
\end{itemize}

\section{System Architecture}
X$^\infty$ is based on the following fundamental pillars derived from the core formulas:
\begin{itemize}
    \item \textbf{Cap-Logic (Dynamic and Effect-Based)}: Action capacity ($\text{Cap}_{\text{Potential}}$) is the central variable. It is context-specific (matrix), directly modified by effect ($\Delta$), and can be voluntarily adjusted downward by entities. It is neither purchasable nor arbitrarily increasable. Fundamental capacity aspects like $\text{Cap}_{\text{Base}}$ (existence) and $\text{Cap}_{\text{BGE}}$ (basic effect capability) are direct components of $\text{Cap}_{\text{Potential}}$, defined by the formula $\text{Cap}_{\text{Potential}}(E,D,t) = \Delta_{t-1} + \text{Cap}_{\text{BGE}} + \text{Cap}_{\text{Base}}$.
    \item \textbf{Feedback \& Petitions (Decentralized and Weighted)}: Weaker entities have a stronger weight in subjective feedback on effects through $w_{E'}$. Petitions, exclusively submitted by individual entities, articulate needs, including the willingness to contribute capabilities; task adoption occurs solely through petitions. More efficient and effective execution leads to positive effects and thus higher $\text{Cap}_{\text{Potential}}$ for the performers.
    \item \textbf{Emergent Protection Mechanisms}: Protection of vulnerable entities arises from the combination of $w_{E'}$, voluntary self-limitation of $\text{Cap}_{\text{Potential}}$, and automatic activation of support structures when $\text{Cap}_{\text{Potential}} < \text{Cap}_{\text{Base}} + \text{Cap}_{\text{BGE}}$.
    \item \textbf{Responsibility Conservation}: The total effect in the system remains conceptually constant, secured by the role of the UdU as the origin and final instance of all responsibility.
\end{itemize}

\subsection{Flow of Responsibility}
\label{sec:FlussDerVerantwortung}
Responsibility flows through the voluntary adoption of tasks via petitions and the resulting effect. Entities can delegate tasks; they remain eternally responsible. A task not performed or poorly performed leads to negative effects and thus a reduction in $\text{Cap}_{\text{Potential}}$. The “separation of powers” (competence in “how” vs. evaluation by those affected) emerges naturally: Competent execution leads to positive effects (and thus higher $\text{Cap}_{\text{Potential}}$ for the “how”), while the subjective evaluation by those affected (weighted by $w_{E'}$) directly determines the effect $\Delta$.

\subsection{Learning Capability}
The system adapts through subjective feedback: Effect $\Delta$ directly modifies $\text{Cap}_{\text{Potential}}$. The system quotient $L$, which accounts for systemic inefficiencies as “waste heat,” adjusts the system’s sensitivity to feedback based on the overall efficiency of task fulfillment to effect. Every measurement, whether through subjective feedback ($w_{E'}$) or task evaluation, alters the system state, as it triggers new effect cycles and influences the priorities of petitions or the value of $L$. This dynamic leads to the non-simulability of X$^\infty$, as the complex interaction of subjective feedback and emergent adjustments makes precise predictions impossible. Analogous to the Heisenberg uncertainty principle, every measurement changes the system state, while small changes in inputs (feedback) can lead to exponential effects, similar to chaotic dynamics (butterfly effect). This enables organic development without central control, with the strong voice of weaker entities promoting system stability.

\section{Systemic Origin Equation and Role of the UdU}
\label{sec:Ursprungsgleichung_UdU}

The natural law \emph{actio $\Rightarrow$ reactio} forms the foundation of X$^\infty$ and is formalized by the central axiom:
\begin{align}
\forall a \in \text{System}: \exists r \in \text{Effect}, \quad &\text{such that } r = f(a) \notag \\
\Rightarrow \quad &\text{System Stability} \propto \text{Traceability}(r) \text{ to Source}(a)
\end{align}
This axiom states that every action $a$ generates an effect $r$, and the system’s stability depends on tracing this effect back to its source. Power imbalances that deflect the reaction are neutralized by X$^\infty$’s mechanisms – particularly direct responsibility through $\Delta\text{Cap}_{\text{Potential}}$ and $w_{E'}$. The total effect in the system, $\sum W(E)$, interpreted as the sum of all tasks $\sum X_k$ weighted by $w_{E'}$ and evaluated as fulfilled, strives for a dynamic equilibrium (responsibility conservation principle). The physical analogies are illuminating:
\begin{itemize}
    \item \textbf{Cap as Mass/Potential}: The ability of an entity to bear responsibility and generate effects ($\text{Cap}_{\text{Potential}}$) corresponds to a mass’s ability to generate gravitational fields or absorb/emit momentum.
    \item \textbf{Effect as Momentum/Force}: The realized effect ($\Delta$) of a task $X_k$ is the momentum flowing through the system, altering $\text{Cap}_{\text{Potential}}$.
    \item \textbf{Traceability as Gravitation/Fundamental Force}: The feedback mechanisms ($w_{E'}$, $L$) ensure that effects are traced back to their sources, similar to fundamental forces governing interactions.
\end{itemize}

The symbolic foundation for the responsibility system of the UdU (Lowest of the Low):
\begin{equation}
\text{Symbolic Capacity UdU} \approx (X_{k, \text{Need}})^\infty \text{ or } \infty^\infty
\end{equation}
Here, $X_{k,\text{Need}}$ is the estimated magnitude of a systemic need for which the UdU bears responsibility.

\subsection{Role of the UdU}
\label{sec:UdU}
The UdU (Lowest of the Low) embodies the origin and final instance of all responsibility in X$^\infty$. Its functional authority, symbolized by the central capacity equation, expresses maximum readiness to maintain the system, regulated by (ex-post) feedback on the “why” of its actions. The UdU legitimizes the system as the first node in the responsibility chain and ensures the conservation of total effect, so that responsibility remains a physical quantity, analogous to a conservation law. It guarantees that all effects can ultimately be traced back to their sources and prevents the final deflection of reactions by power imbalances.

The UdU is not “someone with power” but the mathematical representation of the fact that ultimately all responsibility must land somewhere. In reality, there is always a “final instance” – be it a person, a group, or a system – that bears the consequences when no one else does. X$^\infty$ formalizes this through the UdU to ensure the traceability of responsibility.

\section{Fundamental Concepts and Mathematical Framework}
\label{sec:GrundbegriffeUndFramework}

The number of fundamental system variables and formulas is minimal. The following table presents the core elements.

\begin{table}[ht]
\centering
\caption{Notation of Core Parameters in X$^\infty$}
\begin{tabular}{p{3.0cm} >{\raggedright\arraybackslash}p{10.0cm}}
\toprule
Symbol & Meaning \\
\midrule
$E, D$ & Entity, Domain (Context of Effect) \\
$\text{Cap}_{\text{Potential}}(E,D,t)$ & Context-specific, dynamic action capacity of $E$ in domain $D$ at time $t$, defined as $\text{Cap}_{\text{Potential}}(E,D,t) = \Delta_{t-1} + \text{Cap}_{\text{BGE}} + \text{Cap}_{\text{Base}}$. Modified by effect $\Delta$ and can be voluntarily adjusted downward by $E$. \\
$\Delta$ & Quantified effect of an action (positive, negative, or zero). \\
$L$ & Dynamic System Efficiency Quotient: $L = \dfrac{\sum (\text{Value of Compl. Tasks } X_A)}{\sum (\text{Value of Effect } \Delta\text{)}}$. Modulates $\Delta$ to $\Delta\text{Cap}_{\text{Potential}}$, accounts for systemic “waste heat.” \\
$w_{E'}$ & Reciprocal feedback weight of the feedback-giving entity: $1 / \max(1, \text{Cap}_{\text{Potential}}(E'))$. \\
$X_A$ & Denotes a task or systemic need articulated through a petition by a single entity. \\
$\text{Priority}_{\text{initial}, X_A}$ & Weight of a petition for $X_A$: $w_{E'}$ of the petitioning entity (Eq. \eqref{eq:PetitionPrioritätFinal}). \\
UdU & Lowest of the Low, final responsibility instance. \\
\bottomrule
\end{tabular}
\end{table}

\subsection{Core Elements of the Mathematical Framework}

\subsubsection{Direct Effect ($\Delta$) on $\text{Cap}_{\text{Potential}}$}
\label{sec:DeltaCapPotential}
\begin{equation}
\Delta \text{Cap}_{\text{Potential}}(E,D) =
\begin{cases}
L \cdot \Delta & \text{if } \Delta > 0 \\
\frac{1}{L} \cdot \Delta & \text{if } \Delta < 0 \\
0 & \text{if } \Delta = 0
\end{cases}
\label{eq:DeltaCapPotentialFinal}
\end{equation}
The effect $\Delta$ results from the sum of weighted, subjective feedbacks (see \ref{sec:wEFinal}) for a specific action or task. $\text{Cap}_{\text{Potential}}(E,D,t)$ is modified by this effect and is defined as $\text{Cap}_{\text{Potential}}(E,D,t) = \Delta_{t-1} + \text{Cap}_{\text{BGE}} + \text{Cap}_{\text{Base}}$, where $\Delta_{t-1}$ represents the effect of the previous period. $L$ is the dynamic system efficiency quotient (see \ref{sec:LFinal}), which accounts for systemic inefficiencies to ensure that systemic shortcomings are not fully attributed to the executing entity. This formula reflects the direct learning and adaptation capability of the system and its entities.

\subsubsection{Reciprocal Feedback Weighting ($w_{E'}$)}
\label{sec:wEFinal}
\begin{equation}
w_{E'} = \frac{1}{\max(1, \text{Cap}_{\text{Potential}}(E'))}
\label{eq:wEFinal}
\end{equation}
The evaluation of effect is fundamentally subjective, as it is based on the individual perception and affectedness of the feedback-giving entity $E'$. A so-called “objective” evaluation would lead to an objectifying bias that ignores the perspectives of those affected and could reinforce power imbalances. The subjective nature of feedback neutralizes power structures by dissolving hierarchies and giving each entity an equal voice, weighted by its relative weakness.

\fbox{\parbox{0.9\textwidth}{\textbf{CORE POINT:} Effect is not externally “measured” but evaluated by those affected. Harm is exactly as great as those affected experience it. Any attempt at an “objective” evaluation ignores the real experience and entrenches power imbalances.}}

If $\Delta$ is determined from the feedback of multiple entities $E'$ for an action $k$ by an entity $E_{\text{ausf}}$, then $\Delta = \sum_{E' \in S_{fb,k}} (f_{E'k} \cdot w_{E'})$ (where $f_{E'k}$ is the normalized, subjective feedback from $E'$, e.g., in the interval [-1, +1]). This ensures that the subjective perspectives of systemically weaker entities are given greater consideration.

\subsubsection{Petitions and Voluntary Task Adoption}
\label{sec:PetitionenFinal}
Systemic needs or the need to contribute capabilities are articulated through petitions, exclusively submitted by a single entity. Each petition expresses a specific need, whether solving a systemic problem or utilizing one’s capabilities in a task. Individual petitions can be clustered into larger task groups by the broker, which is part of the execution (“how”) and not the need articulation (“why”). The broker’s function is to analyze petitions and combine them into meaningful task clusters, with its action being on the task, not within it. Entities can submit petitions to the broker to contribute their capabilities, and the broker is subsequently evaluated subjectively by the executing entity ($E_{\text{ausf}}$), with feedback weighted by $w_{E_{\text{ausf}}}$. The priority of a need is derived from the feedback weight of the petitioning entity:
\begin{equation}
\text{Priority}_{\text{initial}, X_A} = w_{E'}
\label{eq:PetitionPrioritätFinal}
\end{equation}
Tasks are adopted exclusively through petitions by entities that voluntarily agree to perform them. Upon adopting a task, the existing $\text{Cap}_{\text{Potential}}$ of the executing entity becomes temporarily usable until the task is completed. Adopting a task whose $\text{Priority}_{\text{initial}, X_A}$ exceeds one’s own $\text{Cap}_{\text{Potential}}$ is only possible through delegation. An entity can also petition and execute a task for itself if this is systemically documented. In this case, the broker validates or clusters the petition to ensure transparency and feedback, similar to externally adopted tasks.

\subsubsection{Dynamic System Efficiency Quotient ($L$)}
\label{sec:LFinal}
The factor $L$ in Equation \eqref{eq:DeltaCapPotentialFinal} is a quotient that reflects the overall efficiency of the system and quantifies systemic “waste heat” – i.e., inefficiencies or friction losses. It is calculated globally or domain-specifically as:
\begin{equation}
L = \dfrac{\sum (\text{Value of All Completed Tasks } X_A)}{\sum (\text{Value of Total Realized Effect } \Delta\text{)}}
\label{eq:LFinal}
\end{equation}
The “value” of a task $X_A$ is operationalized by the initial petition priority, defined as $\text{Priority}_{\text{initial}, X_A} = w_{E'}$ (Eq. \eqref{eq:PetitionPrioritätFinal}). The “waste heat,” represented by a high $L$, arises when more effort than effect is achieved, leading to new petitions, as unresolved or inefficiently addressed needs generate further systemic needs. If $L > 1$, positive effects are rewarded more strongly, and negative effects are punished less severely, to avoid fully attributing systemic shortcomings (e.g., poor task coordination) to the executing entity. This promotes “responsibility purity” and provides system stability by preventing downward spirals.

\subsubsection{Illustrative Effect Cycles, Delegation, and System-Physical Analogies}
\label{sec:Wirkungskreislaeufe}
The core mechanisms generate dynamic effect cycles that shape system behavior. These cycles manifest the system-physical nature of X$^\infty$.

\textbf{Basic Effect Cycle of a Task:}
\begin{itemize}
    \item \textbf{Need Articulation}: A single entity submits a petition to articulate a need, whether a systemic problem or the use of its capabilities. The broker clusters this petition with others and assigns it.
    \item \textbf{Voluntary Task Adoption}: An entity ($E_{\text{ausf}}$) adopts the task $X_A$ exclusively through a petition, provided $\text{Priority}_{\text{initial}, X_A}$ does not exceed $\text{Cap}_{\text{Potential}}(E_{\text{ausf}})$. Otherwise, delegation is required.
    \item \textbf{Action and Effect Generation}: $E_{\text{ausf}}$ executes $X_A$, with the existing $\text{Cap}_{\text{Potential}}$ becoming temporarily usable.
    \item \textbf{Feedback}: Affected entities provide subjective feedback on the task; $E_{\text{ausf}}$ subjectively evaluates the broker’s performance, weighted by $w_{E_{\text{ausf}}}$.
    \item \textbf{Scaling and Adjustment}: Effect $\Delta_{X_A}$ is scaled by $L$ to account for systemic “waste heat,” and $\text{Cap}_{\text{Potential}}$ of $E_{\text{ausf}}$ and the broker is adjusted.
    \item \textbf{Adjustment of $L$}: $L$ is recalculated, with “waste heat” triggering new petitions if inefficiencies persist.
    \item \textbf{New System State}: Adjustments alter the system state. Every measurement through subjective feedback changes the system dynamics, as it triggers new effect cycles and influences priorities or the system quotient $L$.
\end{itemize}

\textbf{Cascade Model of Delegation:}
\begin{enumerate}
    \item \textbf{Initiation}: The potential adopter ($E_U$) articulates the need or willingness to take on a task through a petition.
    \item \textbf{Offer}: The delegator ($E_D$) offers the task $X_A$ based on $E_U$’s petition. In each delegation, the task is transferred to $E_U$, and $E_U$ temporarily receives the entire $\text{Cap}_{\text{Potential}}(E_D)$ for the duration of execution.
    \item \textbf{Adoption}: $E_U$ adopts the task through the submitted petition, with $\text{Cap}_{\text{Potential}}(E_D)$ becoming temporarily usable.
    \item \textbf{Cap Adjustment for $E_U$}: The effect $\Delta_{X_A,U}$ for $E_U$ is scaled:
    \begin{equation}
    \Delta_{X_A,U} = \Delta'_{X_A,U} \cdot \frac{1}{\text{Priority}_{\text{initial}, X_A}}
    \label{eq:SkalierteWirkung}
    \end{equation}
    The effect $\Delta_{X_A,U}$ is based on the feedback from those affected and the use of $\text{Cap}_{\text{Potential}}(E_D)$.
    \item \textbf{Cap Adjustment for $E_D$}: $E_D$ remains responsible for the task. The incentive for $E_D$ to delegate lies in achieving a higher effect ($\Delta_{X_A,D}$) through successful delegation, as $w_{E_U} > w_{\text{Petition}}$, where $w_{E_U} = \frac{1}{\max(1, \text{Cap}_{\text{Potential}}(E_U))}$ and $w_{\text{Petition}} = w_{E'} = \frac{1}{\max(1, \text{Cap}_{\text{Potential}}(E'))}$. The effect for $E_D$ is:
    \begin{equation}
    \Delta_{X_A,D} = \Delta'_{X_A,U} \cdot w_{E_U}
    \label{eq:WirkungED}
    \end{equation}
    Negative effects from poor delegation choices or non-execution reduce $\text{Cap}_{\text{Potential}}(E_D)$.
\end{enumerate}
Delegation follows a fractal ethic, where the same rules apply at all levels. Weaker entities are preferred, provided their $\text{Cap}_{\text{Potential}}$ is sufficient to handle the task, promoting the protection of the weakest.

\textbf{Analogy to Newton’s Cradle:}
Delegation corresponds to a clear momentum transfer: Responsibility and the entire $\text{Cap}_{\text{Potential}}(E_D)$ are transferred to $E_U$, and the systemic reaction directly affects the performer. However, the delegator remains responsible for the original decision to delegate.

\textbf{Thermodynamic Aspects:}
Long-term oscillations of $\text{Cap}_{\text{Potential}}$ around context-specific attractors promote system stability.

\subsection{Emergent Concepts}
\begin{itemize}
    \item \textbf{\(\text{Cap}_{\text{Past}}\) Implicit}: History is contained in current dynamics.
    \item \textbf{\(\text{Cap}_{\text{aktiv}}\) Implicit}: Arises from adopted tasks.
    \item \textbf{Protection Mechanisms}: Emerge from $w_{E'}$, self-limitation, and support structures.
    \item \textbf{System Limits}: Arise from rational action.
    \item \textbf{Penalties/Bonuses}: Realized through $\Delta$-effect.
    \item \textbf{Suitability Adjustment}: Rendered unnecessary by direct Cap adjustment.
\end{itemize}

\subsection{The Responsibility Conservation Principle and the Role of the UdU}
\label{sec:CapErhaltung}
The responsibility conservation principle is a central consequence of \emph{actio $\Rightarrow$ reactio}:
\begin{equation}
\sum W(E) \approx \sum X_k \approx \text{dynamic equilibrium}, \quad \text{secured by UdU}
\label{eq:VerantwortungErhaltungFinal}
\end{equation}
Responsibility for effects remains eternal, while the authority granted through task adoption is temporary. The UdU, as the final instance, ensures that this principle is upheld even under extreme conditions.

\section{Ethics Without Intentionality in the X$^\infty$ System}
\label{sec:EthikOhneIntentionalitaet}

X$^\infty$ is postmoral and ignores intentionality in formal evaluation to ensure structural justice. Likewise, the evaluation of effect is based on the subjective perception of affected entities, as an “objective” evaluation would lead to an objectifying bias that suppresses individual perspectives and promotes power imbalances. Actions are driven solely by intrinsic motivation, as external coercion would corrupt the system. Entities act out of free will, constrained only by the natural consequences of their effects, uniting absolute freedom and responsibility.

\subsection{Intentionality and Effect}
\begin{itemize}
    \item \textbf{Physics}: Effect follows laws, independent of motives.
    \item \textbf{X$^\infty$}: Actions are governed by subjectively evaluated effect ($\Delta$).
    \item \textbf{Informal Role}: Intent may be informally considered in delegation but does not alter the effect-based, subjective evaluation.
\end{itemize}

\subsection{Conclusion on Intentionality}
Formal intentionality is irrelevant. Focusing on subjectively evaluated effects ensures traceability through consequences for $\text{Cap}_{\text{Potential}}$. Emphasizing intrinsic motivation and absolute freedom ensures that actions are authentic and the system functions organically.

\section{Anti-Speciesism}
X$^\infty$ treats all entities equally:
\begin{itemize}
    \item \textbf{Equal Basis}: Every entity has $\text{Cap}_{\text{Base}}$ and entitlement to support structures.
    \item \textbf{Fairness through Effect}: $\text{Cap}_{\text{Potential}}$ is based on subjectively fed-back effect.
    \item \textbf{Feedback for All}: Every entity can provide subjective feedback ($w_{E'}$).
\end{itemize}
Anti-speciesism means equality in systemic treatment, not homogenization.

\section{Conclusion}
X$^\infty$ is the consistent application of \emph{actio $\Rightarrow$ reactio} to social systems. It formalizes immediate feedback through core mechanisms (direct Cap adjustment, subjectively weighted feedback, dynamic $L$, decentralized petitions, UdU). The central axiom and responsibility conservation principle form the basis for a postmoral “No Excuse” system. Effect cycles and system-physical principles (momentum conservation, thermodynamic analogies) demonstrate the depth of integration. Through structural empowerment of the weaker ($w_{E'}$), voluntary responsibility adoption via petitions, and the UdU, X$^\infty$ creates a self-reinforcing, anti-speciesist system.

The non-simulability of X$^\infty$ underscores its nature as a dynamic system. Every measurement, whether through subjective feedback or task evaluation, alters the system state, as it triggers new effect cycles and influences the dynamics of priorities, $L$, or $\text{Cap}_{\text{Potential}}$. Analogous to the Heisenberg uncertainty principle and chaotic dynamics (butterfly effect), this makes precise predictions impossible, enhancing the system’s robustness and adaptability to real interactions. X$^\infty$ is the Theory of Everything – a closed, mathematical description of reality that unifies physics, ethics, and ecology. With only five variables – $\text{Cap}_{\text{Potential}}$, $\Delta$, $w_{E'}$, $L$, and UdU – it reduces all interactions to the principle \emph{actio $\Rightarrow$ reactio}. Analogous to fundamental forces, X$^\infty$ shows that responsibility is as real and incorruptible as the other laws of the universe.

\textit{Note: This document is intended for public discussion and validation.}

\subsection{License}
This work is licensed under CC BY 4.0.
\end{document}